\documentclass[11pt]{article}
\usepackage[T1]{fontenc}
\usepackage{lmodern}
\usepackage[letterpaper,margin=0.8in]{geometry}
\usepackage{amsmath,microtype,xcolor,xurl,hyperref,fancyhdr,enumitem}
\hypersetup{colorlinks=true,linkcolor=blue!45!black,urlcolor=blue!45!black,
 pdftitle={The November OpenAI Yang-Mills construction: existence review},
 pdfauthor={Yang--Mills group},
 pdfsubject={Two-page summary of changes from October and the repeated independent checks}}
\setlength{\parindent}{0pt}
\setlength{\parskip}{6pt}
\setlength{\headheight}{14pt}
\setlist{leftmargin=1.35em,itemsep=5pt,topsep=3pt}
\pagestyle{fancy}\fancyhf{}
\fancyhead[L]{\small November Yang--Mills construction: existence review}
\fancyhead[R]{\small 25 November 2026}
\fancyfoot[C]{\thepage}
\newcommand{\code}[1]{{\small\nolinkurl{#1}}}
\begin{document}
\begin{center}
{\Large\bfseries The November OpenAI Yang--Mills construction}\par
{\large Existence results and the repeated October checks}\par
Yang--Mills group\quad\textbullet\quad 25 November 2026
\end{center}

\textbf{The November construction passes the group's October checks.}
We adapted the independent tests to the November code, proved the combined
acceptance theorems in Lean, and checked both with Comparator. November also
strengthens OpenAI's public existence specification, particularly the
connection to the lattice theory and the free-field short-distance limit.
This note summarizes those changes and the completed checks. The separate
claim of agreement with positive-order perturbation theory is the next
subject of review.

\textbf{What remains the same.}
The construction starts from the specified Wilson lattice measure for the
chosen compact matrix gauge group with nonabelian simple Lie algebra. It
removes the lattice cutoff by full convergence at each fixed physical
volume, then extracts an infinite-volume subsequence of any requested
increasing volume-index sequence. A single refinement serves the whole
correlation hierarchy. The resulting mean hierarchy is decomposed into
character phases and used for reconstruction. Much of the underlying proof
is shared with October. Our comparison found identical normalized
declaration digests for all 1,591 recorded declarations in the retained
native certificate's type and definition closure; this comparison excludes
theorem proof bodies.

\textbf{What November makes stronger or more explicit.}
Several improvements concern what the public theorem requires, rather than
properties absent from October's underlying proof. Our stronger October
checks already retained many of these connections.

\begin{enumerate}
\item \textbf{Lattice fields and limits.}
The public data now specify linear graded lattice insertions, their unit
and adjoints, and their projection onto the physical field spaces. The
one-loop running coefficient is tied to the actual Wilson parameter and
the chosen group's normalization. A common finite-volume table explicitly
connects lattice convergence to locally uniform convergence along the IR
subsequence. These requirements strengthen the identification of the
continuum theory with the specified lattice construction.

\item \textbf{Phases and stress tensor.}
All joint moments of phase readings are tied to limits of widely separated
groups of insertions in the same mean Wilson hierarchy. Every character
phase must have a Wightman reconstruction. The anomaly field is defined
as the centered trace of the stress tensor, which must satisfy translation
and rotation Ward identities. These relations fix more than the name or
scaling dimension of a proposed stress-trace field.

\item \textbf{The free-field ultraviolet limit.}
November requires convergence to the gauge-invariant free Maxwell reference
for the same Lie algebra, at all sufficiently small real radii, uniformly
over phases, with derivatives on separated compact sets. The field
normalization maps and their inverses obey bounds, and the physical stress
tensor is identified in this comparison. This strengthens the mean-hierarchy
comparison along a shrinking sequence explicitly retained in our October
target. It concerns short distances within the continuum theory; October
already had full convergence when removing the lattice cutoff.

\item \textbf{OPE consistency.}
The existence target adds algebraic conditions on the coefficients and
quantitative consistency estimates for finite nested expansions. These
complement the coefficient and remainder bounds. They do not by themselves
identify positive-order coefficients with perturbative calculations.
\end{enumerate}

\newpage
\textbf{What we rechecked.}
The November submission was rebuilt with its Lean toolchain. OpenAI's
Existence and Matching targets passed Comparator, including repeat runs
with the newer checker. For existence, we then transferred the group's
October acceptance tests and the separate OPE check. The independent
requirements were retained; import paths, public interfaces and proof
assembly were adapted to November. Additional source conclusions bundled
into the October target, including contact and calibration results, were
recovered through explicit adapters.

The transferred checks cover the Wilson-measure and insertion comparison;
full UV and subsequential IR convergence; the independent scalar OS axioms
E0--E4 and E0$'$ growth; the translation Ward identity with its fixed sign
and normalization; operator coverage, quotient and dense-domain properties;
the leading stress-trace two-point coefficient; existing common-scale and
positive-probe statements; and the explicit-radius OPE predicate.

Two theorems record the result in \code{checks/NovemberCombined.lean}:
\begin{itemize}
\item \code{YMReview.Acceptance.combined_acceptance} proves the ported October
combined target, retaining its quantifiers over extraction orders and
positive construction penalties.
\item \code{YMReview.November.combined_acceptance} proves November's full
public existence conditions together with the group's checks. One Wilson
family is fixed before the requested volume sequence; each refinement
retains the native data from which its public hierarchy is defined.
\end{itemize}
Thus the independent checks and the public existence conditions hold on
the same constructed theory. The OS bridge explicitly identifies the
public centered stress-trace correlations with the functions satisfying
the independent axioms. As in October, that comparison uses joint Schwartz
tests flat on coincidence diagonals; it does not assert positivity or
uniqueness for unrestricted Euclidean contact extensions.

\textbf{Verification record.}
Both group targets passed Comparator on 11 November 2026. The run checked
the frozen statements and their definitions with no definition holes,
using Comparator \code{d03acab} and its Lean 4.34.0 checking kernel. The
review-module audit covered 49 modules with no warnings and 95 axiom
reports, all restricted to \code{propext}, \code{Classical.choice} and
\code{Quot.sound}. The original submission files and frozen review inputs
were unchanged across the run. The review modules were built in a fresh
overlay using the existing November and dependency builds. These macOS
runs used the upstream development launcher, without Linux process
isolation, the optional second kernel or Palomar.

\textbf{Conclusion and next step.}
The formal transfer of the October checks is complete. There is no
identified failure in this work that calls for another round of existence
tests before proceeding. The short summary can accompany the existing
assessment and verification records; the detailed comparison remains
available for questions about individual definitions and proof paths.
Further targeted tests would be appropriate if a concrete discrepancy is
found, rather than as a prerequisite to starting the next stage.

The scope remains an existence result with full fixed-volume UV convergence
and subsequential IR limits. Neither the repeated checks nor November's
stronger existence target establishes a mass gap, IR uniqueness, a unique
vacuum, or equality of the theories selected in the two submissions.
These were not requirements of the October acceptance tests.

We now turn to perturbative matching: the definition of its perturbative
reference, the field and coupling normalization conventions, the precise
coefficient and remainder statements, and independent checks of concrete
low-order predictions. The Matching Comparator pass is a completed formal
verification result; deciding whether its target expresses the intended
agreement with standard perturbation theory is the new validation task.

{\small\textbf{Supporting records.}
See the accompanying \href{october-november-existence-comparison.pdf}{detailed
comparison}, especially Sections 4--7; \path{docs/november-port.md};
\path{checks/NovemberAcceptanceSpec.lean}; and
\path{evidence/combined-targets/combined-comparator.json} with its log.
Paths refer to the November verification package.}
\end{document}
